\documentclass[11pt,a4paper]{article}
% Estilo común. Editar aquí para mantener coherentes los tres documentos.
% Adaptación de las plantillas del Seminario Agentes Inteligentes y LLM (UNLu).
\usepackage[utf8]{inputenc}
\usepackage[T1]{fontenc}
\usepackage[spanish,es-nodecimaldot]{babel}
\usepackage[default]{sourcesans}
\usepackage{microtype}
\usepackage[a4paper,left=23mm,right=23mm,top=23mm,bottom=23mm,headheight=15pt,headsep=8mm,footskip=11mm]{geometry}
\usepackage{xcolor}
\usepackage{graphicx}
\usepackage{enumitem}
\usepackage{titlesec}
\usepackage{fancyhdr}
\usepackage{needspace}
\usepackage[skins,breakable]{tcolorbox}
\usepackage{hyperref}
\usepackage{xurl}

\definecolor{azul}{HTML}{1A465E}
\definecolor{petroleo}{HTML}{2F859A}
\definecolor{texto}{HTML}{294D60}
\definecolor{suave}{HTML}{EFF7F9}
\definecolor{borde}{HTML}{CFE2E8}
\definecolor{gris}{HTML}{586D79}

\hypersetup{colorlinks=true,urlcolor=petroleo,linkcolor=azul,citecolor=azul,
  pdfauthor={Aprendizaje Automático - Universidad Nacional de Luján}}
\urlstyle{same}
\color{texto}
\setlength{\parindent}{0pt}
\setlength{\parskip}{6pt}
\linespread{1.10}
\setlength{\emergencystretch}{1.5em}
\clubpenalty=10000
\widowpenalty=10000

% Más aire antes del título que después: el título acompaña su contenido.
\titleformat{\section}{\Large\bfseries\color{azul}}{}{0pt}{}
\titlespacing*{\section}{0pt}{24pt}{0pt}
\titleformat{\subsection}{\large\bfseries\color{azul}}{}{0pt}{}
\titlespacing*{\subsection}{0pt}{18pt}{0pt}
\setlist[itemize]{leftmargin=16pt,itemsep=4pt,topsep=4pt,parsep=0pt,label=\textcolor{petroleo}{\textbullet}}
\setlist[enumerate]{leftmargin=20pt,itemsep=6pt,topsep=4pt,parsep=0pt,label=\textcolor{petroleo}{\bfseries\arabic*.}}

\pagestyle{fancy}
\fancyhf{}
\fancyhead[L]{\small\color{gris}Universidad Nacional de Luján}
\fancyhead[R]{\small\color{gris}Aprendizaje Automático · 2026}
\fancyfoot[L]{\small\color{gris}\pieDocumento}
\fancyfoot[R]{\small\color{petroleo}\thepage}
\renewcommand{\headrulewidth}{0.35pt}
\renewcommand{\footrulewidth}{0pt}
\newcommand{\pieDocumento}{Trabajo práctico integrador}

\newcommand{\cabecera}[3]{%
  {\small\bfseries\color{petroleo}\MakeUppercase{#1}\par}
  \vspace{4pt}
  {\fontsize{24}{28}\selectfont\bfseries\color{azul}#2\par}
  \vspace{7pt}
  {\large\color{gris}#3\par}
  \vspace{6pt}}

\newtcolorbox{destacado}[1]{enhanced,breakable,colback=suave,colframe=borde,
  boxrule=0.5pt,arc=2mm,left=12pt,right=12pt,top=9pt,bottom=9pt,
  borderline west={2pt}{0pt}{petroleo},before skip=12pt,after skip=10pt,
  title={#1},colbacktitle=suave,coltitle=azul,fonttitle=\bfseries,
  attach title to upper={\par\smallskip},fontupper=\normalsize}

\newcommand{\orientacion}[1]{{\color{gris}#1}\par}
\newcommand{\campo}[1]{\textcolor{petroleo}{\textit{[#1]}}}
% Encabezado compacto para los ejemplos consecutivos de la consigna.
\newcommand{\lineadesarrollo}[1]{\par\vspace{12pt}\noindent{\large\bfseries\color{azul}#1\par}\nopagebreak\vspace{-2pt}}

\hypersetup{pdftitle={Trabajo práctico integrador - Aprendizaje Automático 2026}}

\begin{document}
\cabecera{Trabajo práctico integrador}{De una situación real\ a una solución de IA}{Consigna y orientaciones para el desarrollo}

\begin{destacado}{La propuesta}
Diseñen e implementen una solución de inteligencia artificial para abordar una situación concreta, utilizando \textbf{datos reales}. El resultado será una prueba de concepto reproducible, acompañada por una explicación de sus decisiones, resultados y posibles usos.
\end{destacado}

\section*{Qué buscamos con este trabajo}
El trabajo propone conectar los conceptos del curso con una necesidad reconocible: una pregunta, una tarea o una decisión que podría beneficiarse de un modelo o sistema de IA. Interesa que puedan pasar del problema a su representación, elegir un enfoque pertinente y analizar qué aporta la solución obtenida.

El alcance debe ser realizable con los recursos disponibles. Una notebook bien construida puede ser suficiente, pero también pueden desarrollar scripts, una aplicación u otra forma de prototipo. La complejidad de la interfaz o la cantidad de técnicas utilizadas no reemplazan la calidad del planteo y de la evaluación.

\section*{Componentes del trabajo}
\begin{enumerate}
\item \textbf{Situación y propósito.} Describan el contexto, quién podría utilizar la solución y qué necesidad atendería. Delimiten una tarea concreta y expliquen qué resultado sería útil para esa persona u organización.
\item \textbf{Datos y representación.} Identifiquen la fuente, las características y las limitaciones de los datos. Expliquen cómo los transforman en una representación adecuada: atributos, imágenes, textos, estados, acciones o hechos y reglas, según corresponda.
\item \textbf{Enfoque, entrenamiento e implementación.} Justifiquen la técnica y construyan una prueba de concepto. Cuando corresponda, entrenen o ajusten el modelo y documenten los datos, la configuración y el procedimiento seguido. Si usan un modelo preentrenado sin ajustarlo, justifiquen esa decisión.
\item \textbf{Evidencia y uso.} Evalúen el resultado con criterios apropiados. Analicen errores y limitaciones, e imaginen cómo se integraría el modelo o sistema en la situación planteada.
\end{enumerate}

\begin{destacado}{Datos reales, con procedencia identificable}
Pueden usar fuentes públicas, institucionales, científicas o un relevamiento propio. Deben indicar de dónde provienen los datos y cómo acceder a ellos. \textbf{No se admitirán datos sintéticos ni generados artificialmente.} Si utilizan conocimiento simbólico, identifiquen también la procedencia de los hechos y reglas del dominio.
\end{destacado}

\newpage
\cabecera{Elección del problema}{Líneas posibles de desarrollo}{Orientaciones para encontrar un alcance propio}

Las siguientes líneas son ejemplos, no una lista cerrada. Elijan una situación que puedan explicar y un conjunto de datos que permita trabajar sobre ella. No es necesario utilizar todos los contenidos del curso ni combinar enfoques si esa combinación no aporta al problema.

\lineadesarrollo{Aprendizaje sobre datos estructurados}
Abordar una tarea de clasificación, predicción o análisis a partir de registros reales. Por ejemplo, estimar la demanda de un servicio con datos históricos, clasificar eventos registrados o reconocer patrones en un conjunto de mediciones. La pregunta debe permitir interpretar para qué serviría la predicción y qué consecuencias tendría equivocarse.

\lineadesarrollo{Visión por computadora}
Analizar imágenes de un dominio concreto: residuos, especies, productos o defectos registrados, entre otras posibilidades. Pueden entrenar una red o adaptar un modelo preentrenado. Expliquen qué información contienen las imágenes, cómo se obtuvieron las etiquetas y en qué condiciones podría fallar el modelo.

\lineadesarrollo{Procesamiento del lenguaje natural}
Clasificar, recuperar o analizar textos, voz o audio provenientes de una fuente identificada. Por ejemplo, organizar consultas de soporte, reconocer temas en documentos o analizar opiniones. También pueden explorar generación de lenguaje, siempre que definan qué esperan de la salida y cómo comprobarán su utilidad, sus errores y su relación con la información de origen.

\lineadesarrollo{Conocimiento, agentes y búsqueda}
Modelar una situación real mediante estados y acciones, o representar conocimiento explícito para responder consultas y obtener conclusiones. Por ejemplo, planificar una asignación con restricciones y datos operativos, o construir una base de conocimiento a partir de un procedimiento documentado. La calidad de la representación y la justificación de los costos, reglas o criterios de decisión son parte central del trabajo.

\lineadesarrollo{Decisiones y aprendizaje por refuerzo}
Estudiar una decisión secuencial vinculada con un dominio real, utilizando registros o parámetros documentados. Expliciten qué observa el agente, qué acciones puede realizar y qué criterio de recompensa orienta su comportamiento. Justifiquen qué aspectos de la situación representa el entorno de prueba y qué limitaciones introduce esa representación.

\newpage
\cabecera{Desarrollo y evaluación}{Cómo sostener la solución}{Decisiones justificadas y evidencia comprensible}

\section*{Un recorrido de trabajo posible}
Comiencen con una tarea acotada y una primera solución ejecutable de principio a fin. En enfoques que requieren aprendizaje, preparen los datos y entrenen o ajusten el modelo. Documenten la arquitectura o técnica, los hiperparámetros relevantes y el criterio de selección de la configuración final; analicen la evolución de la pérdida o las métricas cuando corresponda.

Cuando corresponda, separen los datos de entrenamiento, validación y prueba y eviten que la información de prueba intervenga en el ajuste. En problemas temporales, consideren el orden de los registros. En sistemas basados en conocimiento o búsqueda, preparen casos reales documentados que permitan observar tanto el comportamiento esperado como las limitaciones.

\section*{Qué evidencia presentar}
La evaluación debe responder si la solución ayuda con la tarea planteada. Elijan medidas y ejemplos que se puedan interpretar en ese contexto:
\begin{itemize}
\item \textbf{Clasificación o predicción:} métricas pertinentes, errores y ejemplos de aciertos y fallos. Una matriz de confusión o una comparación con una referencia sencilla puede ayudar a explicar el resultado.
\item \textbf{Búsqueda, reglas o decisiones:} casos de prueba, soluciones obtenidas, cumplimiento de restricciones, costos o conclusiones producidas y su explicación.
\item \textbf{Lenguaje o generación:} ejemplos analizados con criterios explícitos, incluyendo salidas incorrectas, información no sustentada o respuestas poco útiles.
\end{itemize}
Comparar con una alternativa simple, cuando sea pertinente, ayuda a reconocer qué aporta el enfoque elegido. No se espera un resultado perfecto: también importa identificar por qué falla, en qué situaciones y qué podrían mejorar.

\section*{Qué se tendrá en cuenta}
Se valorarán la claridad del problema, la procedencia y adecuación de los datos, la relación entre el enfoque y la tarea, el funcionamiento reproducible del prototipo y la calidad del análisis. Las conclusiones deberán apoyarse en la evidencia presentada. Consideren además posibles sesgos, límites de uso, costos y consecuencias de los errores que resulten relevantes para su caso.

Expliquen también cómo imaginan la intervención de la solución: \textbf{quién aporta la entrada, qué devuelve el sistema y qué decisión o acción habilita}. Señalen qué debería revisar una persona y qué faltaría para utilizarlo en el contexto real. No es obligatorio desplegar un servicio de producción.

\newpage
\cabecera{Entrega y documentación}{Un repositorio que se pueda recorrer}{Código, datos identificables e instrucciones de ejecución}

\begin{destacado}{Dos instancias en el campus}
\textbf{Propuesta opcional:} hasta 200 palabras con el problema, los datos, el enfoque previsto y el resultado esperado. Permite recibir una devolución previa, no tiene una fecha de entrega obligatoria y no es requisito para presentar el trabajo final.

\textbf{Entrega final obligatoria:} enlace a un repositorio de GitHub o plataforma equivalente con la solución, un README y un \textbf{informe en PDF}, accesible para su revisión, hasta el \textbf{2 de diciembre de 2026 a las 23:59 h}.
\end{destacado}

Incluyan el código, las dependencias y las instrucciones para ejecutar la solución. Si los datos son voluminosos o no pueden redistribuirse, indiquen la fuente, la versión y los pasos de acceso. Eviten publicar credenciales o información personal que no corresponda difundir.

\section*{El README como puerta de entrada}
Presenten brevemente el problema y la solución, identifiquen los archivos principales y enlacen el informe. Detallen el entorno, las dependencias, los datos de entrada y el orden de ejecución, de modo que otra persona pueda reproducir el trabajo. El análisis y la fundamentación se desarrollan en el informe; no hace falta duplicarlos completos en el README.

\section*{Informe y plantilla LaTeX}
El proyecto debe estar acompañado por un \textbf{informe en PDF}, incluido en el repositorio y enlazado desde el README. Desarrollen allí el problema, los datos, el enfoque, la implementación, los resultados, las limitaciones y el uso previsto de la solución. No es necesario copiar el código completo.

Se ofrecen plantillas LaTeX para la propuesta y el informe, adaptadas del Seminario Agentes Inteligentes y LLM. Pueden editarse en Overleaf sin instalar LaTeX. \textbf{El informe es obligatorio; el uso de LaTeX es opcional.} También pueden usar otro editor y exportar a PDF.

\section*{Declaración de uso de IA}
Incluyan en el informe un \textit{disclaimer} o declaración de uso de IA: una nota que indique qué herramientas o modelos utilizaron como asistencia, para qué tareas, qué aportes incorporaron y cómo los verificaron. No es una exención de responsabilidad: deben comprender y poder explicar lo entregado. Si no utilizaron IA como apoyo, indíquenlo expresamente. Los modelos que integran la solución se describen, además, en la metodología.

\textbf{Antes de entregar,} comprueben el acceso al enlace y la ejecución de la solución. Identifiquen las fuentes de datos, modelos, código y documentación utilizados.

\end{document}
